% portada_algebra_lineal2.tex ─────────────────────────────────────────────────
% Portada Álgebra Lineal 2 — Universidad de Guanajuato
% Gráfica: forma canónica de Jordan J = diag(J_1(λ1), J_2(λ2), J_3(λ3))
%   — directamente relacionada con los temas centrales de Álgebra Lineal 2:
%     · Diagonalización y su generalización cuando faltan eigenvectores
%     · Eigenvalores con multiplicidad algebraica mayor a la geométrica
%     · Bloques de Jordan (operadores nilpotentes sobre cada eigenespacio)
%
% ─── EDITAR AQUÍ ─────────────────────────────────────────────────────────────
\newcommand{\numtarea}{1}
\newcommand{\nomalumno}{Ricardo Le\'on Mart\'inez}
\newcommand{\fecha}{21/09/2026}
\newcommand{\nomprofesor}{Herbert Kanarek Blando}
% ─────────────────────────────────────────────────────────────────────────────

\documentclass[tikz, border=0pt]{standalone}

\usepackage[T1]{fontenc}
\usepackage[utf8]{inputenc}
\usepackage[spanish]{babel}
\usepackage{ebgaramond}
\usepackage[default]{sourcesanspro}
\usepackage{xcolor}
\usepackage{amsmath}
\usetikzlibrary{matrix,calc}

\definecolor{ugblue}{HTML}{002D72}
\definecolor{uggold}{HTML}{B8952A}

\begin{document}
\begin{tikzpicture}

  \useasboundingbox (0,0) rectangle (210mm,297mm);
  \fill[white] (0,0) rectangle (210mm,297mm);

  %% ── Barras de color ─────────────────────────────────────────────────────────
  \fill[ugblue] (0,287mm) rectangle (210mm,297mm);
  \fill[uggold] (0,284mm) rectangle (210mm,287mm);
  \fill[uggold] (0,11mm)  rectangle (210mm,14mm);
  \fill[ugblue] (0,0)     rectangle (210mm,11mm);

  %% ── Encabezado institucional ────────────────────────────────────────────────
  \node[font=\fontsize{15.5}{19}\selectfont\sffamily\bfseries\color{ugblue},
        anchor=center] at (105mm,275mm)
    {\MakeUppercase{Universidad de Guanajuato}};
  \node[font=\fontsize{10.5}{14}\selectfont\rmfamily\itshape\color{gray!65},
        anchor=center] at (105mm,267.5mm)
    {Divisi\'on de Ciencias Naturales y Exactas};
  \node[font=\fontsize{10.5}{14}\selectfont\rmfamily\itshape\color{gray!65},
        anchor=center] at (105mm,261mm)
    {Departamento de Matem\'aticas};

  %% ── Divisor superior ────────────────────────────────────────────────────────
  \draw[ugblue!22, line width=0.4pt] (30mm,254.5mm) -- (101mm,254.5mm);
  \fill[uggold]
    (103mm,254.5mm)--(105mm,256.5mm)--(107mm,254.5mm)--(105mm,252.5mm)--cycle;
  \draw[ugblue!22, line width=0.4pt] (109mm,254.5mm) -- (180mm,254.5mm);

  %% ── Forma canónica de Jordan ────────────────────────────────────────────────
  %% J = diag(J_1(λ1), J_2(λ2), J_3(λ3)): bloques triangulares superiores con
  %% el eigenvalor en la diagonal y 1's en la superdiagonal; cero en el resto.
  %% Rejilla: celda 14mm, separación 2.5mm ⇒ paso 16.5mm; matriz 6×6 centrada
  %% en (105mm,186mm). Los realces de bloque se dibujan primero (coordenadas
  %% calculadas a partir de la rejilla) para que queden detrás de los números.
  \fill[ugblue!16,  rounded corners=1.5mm] (55.75mm,186.25mm) rectangle (104.75mm,235.25mm);
  \fill[uggold!28,  rounded corners=1.5mm] (105.25mm,153.25mm) rectangle (137.75mm,185.75mm);
  \fill[ugblue!16,  rounded corners=1.5mm] (138.25mm,136.75mm) rectangle (154.25mm,152.75mm);

  \matrix (M) [
    matrix of math nodes,
    column sep = 2.5mm,
    row sep    = 2.5mm,
    nodes = {
      anchor         = center,
      minimum width  = 14mm,
      minimum height = 14mm,
      font           = \fontsize{17}{17}\selectfont,
    },
  ] at (105mm,186mm) {
    \color{ugblue}\lambda_1 & \color{black!60}1       & \color{black!25}0       & \color{black!25}0       & \color{black!25}0       & \color{black!25}0 \\
    \color{black!25}0       & \color{ugblue}\lambda_1 & \color{black!60}1       & \color{black!25}0       & \color{black!25}0       & \color{black!25}0 \\
    \color{black!25}0       & \color{black!25}0       & \color{ugblue}\lambda_1 & \color{black!25}0       & \color{black!25}0       & \color{black!25}0 \\
    \color{black!25}0       & \color{black!25}0       & \color{black!25}0       & \color{uggold!80!black}\lambda_2 & \color{black!60}1       & \color{black!25}0 \\
    \color{black!25}0       & \color{black!25}0       & \color{black!25}0       & \color{black!25}0       & \color{uggold!80!black}\lambda_2 & \color{black!25}0 \\
    \color{black!25}0       & \color{black!25}0       & \color{black!25}0       & \color{black!25}0       & \color{black!25}0       & \color{ugblue}\lambda_3 \\
  };

  %% Corchetes de la matriz
  \coordinate (nw) at ($(M.north west)+(-3mm, 2mm)$);
  \coordinate (sw) at ($(M.south west)+(-3mm,-2mm)$);
  \coordinate (ne) at ($(M.north east)+( 3mm, 2mm)$);
  \coordinate (se) at ($(M.south east)+( 3mm,-2mm)$);

  \draw[ugblue, line width=1.1pt]
    ($(nw)+(4mm,0)$) -- (nw) -- (sw) -- ($(sw)+(4mm,0)$);
  \draw[ugblue, line width=1.1pt]
    ($(ne)+(-4mm,0)$) -- (ne) -- (se) -- ($(se)+(-4mm,0)$);

  %% ── Divisor inferior ────────────────────────────────────────────────────────
  \draw[ugblue!22, line width=0.4pt] (30mm,121mm) -- (101mm,121mm);
  \fill[uggold]
    (103mm,121mm)--(105mm,123mm)--(107mm,121mm)--(105mm,119mm)--cycle;
  \draw[ugblue!22, line width=0.4pt] (109mm,121mm) -- (180mm,121mm);

  %% ── Título ──────────────────────────────────────────────────────────────────
  \node[font=\fontsize{9}{11}\selectfont\sffamily\fontseries{l}\selectfont
        \color{gray!48}, anchor=center] at (105mm,114mm)
    {\MakeUppercase{Licenciatura en Matem\'aticas}};

  \node[font=\fontsize{42}{48}\selectfont\sffamily\bfseries\color{ugblue},
        anchor=center] at (105mm,97mm)
    {\'Algebra Lineal~2};

  \node[font=\fontsize{27}{31}\selectfont\rmfamily\itshape\color{uggold},
        anchor=center] at (105mm,81mm)
    {Tarea~\numtarea};

  %% ── Separador y datos ───────────────────────────────────────────────────────
  \draw[ugblue!13, line width=0.35pt] (30mm,70.5mm) -- (180mm,70.5mm);

  \node[anchor=west,
    font=\fontsize{7.5}{9}\selectfont\sffamily\color{gray!52}]
    at (30mm,63mm) {\MakeUppercase{Alumno}};
  \node[anchor=west,
    font=\fontsize{7.5}{9}\selectfont\sffamily\color{gray!52}]
    at (120mm,63mm) {\MakeUppercase{Fecha}};
  \node[anchor=west,
    font=\fontsize{15}{18}\selectfont\rmfamily\color{black!85}]
    at (30mm,55.5mm) {\nomalumno};
  \node[anchor=west,
    font=\fontsize{15}{18}\selectfont\rmfamily\color{black!85}]
    at (120mm,55.5mm) {\fecha};

  \draw[ugblue!10, line width=0.3pt] (30mm,48mm) -- (180mm,48mm);

  \node[anchor=west,
    font=\fontsize{7.5}{9}\selectfont\sffamily\color{gray!52}]
    at (30mm,42mm) {\MakeUppercase{Profesor}};
  \node[anchor=west,
    font=\fontsize{15}{18}\selectfont\rmfamily\color{black!85}]
    at (30mm,34.5mm) {\nomprofesor};

  \node[anchor=west,
    font=\fontsize{7.5}{9}\selectfont\sffamily\color{gray!52}]
    at (120mm,42mm) {\MakeUppercase{Calificaci\'on}};
  \draw[ugblue!40, line width=0.5pt, rounded corners=1.2pt]
    (120mm,26.5mm) rectangle (155mm,38.5mm);

  \node[font=\fontsize{7.5}{9}\selectfont\sffamily\color{gray!36},
        anchor=center] at (105mm,21mm)
    {\MakeUppercase{Universidad de Guanajuato~·~Campus Guanajuato}};

\end{tikzpicture}
\end{document}
